\section{Repository difficulties and a concrete repair in Diophantine geometry}
\label{sec:repository-repair}

\subsection{What the inspected packets actually leave open}

The audit in this section refers throughout to ProveIt commit
\path{27a07921176bccbf630084b959700605783ff4aa}.
The index under \path{Logic/PrimitiveRecursiveArithmetic/docs/proofs/}
lists 47 delivered packets and explicitly says that their placement did
not rerun, compile, verify, or edit them.  A packet's own claim of a
compiled native Lean theorem must therefore be distinguished both from an
independent build in this project and from an object-language iPRA
derivation.  We use the packets as inspectable mathematical arguments and
interfaces, rather than treating their index as a correctness certificate
\cite{ProveIt}.

Several obstacles in the packets are already identified with considerable
precision.  It would not advance the program to rename their existing
solutions.  In particular, the Mordell--Weil packet already has contracting
height descent, logarithmic fuel, explicit coordinates, and a method of
obtaining a bounded selector from actual covering representatives.  The
Siegel packet already has a finite-cover-to-cardinality lemma and observes
that covering subspaces of a supplied tuple can be replaced by spans of
subtuples.  Neither result is claimed here as a new repair.

\begin{center}
\small
\begin{tabular}{p{0.19\linewidth}p{0.36\linewidth}p{0.36\linewidth}}
\hline
Packet & Identified missing information & Consequence for this framework\\
\hline
\texttt{puiseux\_ipra} & Arbitrary fast Cauchy names do not decide the
exact coefficient branches needed for full Puiseux output. & Preserve the
exact algebraic input interface; approximation alone does not supply
its zero tests.\\[3pt]
\texttt{mordell\_weil\_ipra} & A finite list of actual rational
representatives for the quotient is still required. & The existing descent
algorithm can consume the list, but does not manufacture it from Selmer
information.\\[3pt]
\texttt{faltings\_ipra} and \texttt{siegel\_ipra} & General positive
height or complete-list bounds are not supplied. & A count bound cannot
serve as a search cutoff.  We prove a useful count ingredient below,
without claiming a height bound.\\[3pt]
\texttt{diophantine\_ipra} & The general effective Roth cutoff is not
proved; quadratic and Liouville-range cases are treated. & Keep the
denominator cutoff distinct from the number of exceptional approximants.\\[3pt]
\path{prime_distribution_ipra} & Chebotarev needs certified
field arithmetic and quantitative analytic estimates; full SW and BV
need the indicated uniform witnesses. & Finite prime enumeration and
finite-precision logarithms do not provide those estimates.\\
\hline
\end{tabular}
\end{center}

These descriptions are claims about the inspected reports, not an
assertion that every stronger theorem is unprovable in iPRA.  Their exact
source paths and blob identifiers are recorded in the accompanying source
manifest.  The following repair targets a specific obligation in
\path{siegel_ipra/paper/05_general_route.tex}: the bounded number of curve
points lying over a proper subspace.  That report sketches a divisor
argument but leaves its arithmetic realization open.  For odd-degree
hyperelliptic equations, the realization can be given directly by one
polynomial norm, with an explicit finite proof.

\subsection{A zero-count theorem from a polynomial norm}

All polynomials and coordinates in this subsection are rational and
finitely encoded.  Rational equality and order mean exact comparisons
after cross multiplication; they are not comparisons of Cauchy names.
Let $f\in\mathbb Q[X]$ have positive odd degree $d$, with nonzero leading
coefficient.  Let
\[
             C_f=\{(x,y)\in\mathbb Q^2:y^2=f(x)\}.
\]
No smoothness assumption is needed for the arithmetic theorem.

For $A,B\in\mathbb Q[X]$, not both zero, define
\[
 G(x,y)=A(x)+B(x)y,\qquad
 N(X)=A(X)^2-f(X)B(X)^2.
\]
In degree inequalities, a zero polynomial simply contributes no term:
thus $2\deg A\le L$ is required only when $A\ne0$, and
$d+2\deg B\le L$ only when $B\ne0$.

\begin{theorem}[Arithmetic fiber bound]
\label{thm:hyperelliptic-fiber}
The polynomial $N$ is nonzero.  Every finite list of pairwise distinct
points of $C_f$ satisfying $G(x,y)=0$ has length at most $\deg N$.
Consequently it has length at most $L$ whenever
\[
        2\deg A\le L,\qquad d+2\deg B\le L
\]
for the nonzero terms.  The theorem has a proof using rational polynomial
arithmetic, bounded finite-list operations, and quantifier-free induction.
\end{theorem}

\begin{proof}
First suppose $B=0$.  Then $A\ne0$ and $N=A^2\ne0$.  If $A=0$ and
$B\ne0$, then $N=-fB^2\ne0$.  If both are nonzero, the degrees of $A^2$
and $fB^2$ have different parity: the former is even and the latter odd.
Their highest terms therefore cannot cancel.  This proves nonvanishing,
and also
\[
 \deg N\le \max\{2\deg A,d+2\deg B\},
\]
with absent terms omitted.

Now take a finite pairwise-distinct list of common zeros.  Group it by
its $x$-coordinate using exact rational equality.  For each different
$x$-coordinate $a$, write $m_a$ for the number of points in that group.
The equation $y^2=f(a)$ has at most two distinct rational solutions:
if $u^2=v^2$, then $(u-v)(u+v)=0$, and the rational field has no zero
divisors.  Thus $m_a\le2$.

Every nonempty group gives $N(a)=0$, by substituting
$A(a)=-B(a)y$ and $y^2=f(a)$.  Hence $X-a$ divides $N$.
If $m_a=2$, let the two different ordinates be $u,v$.  Subtracting the
two equations $A(a)+B(a)u=0$ and $A(a)+B(a)v=0$ gives
$B(a)(u-v)=0$.  Since $u-v\ne0$, we have $B(a)=0$, and then $A(a)=0$.
The factor theorem yields $A=(X-a)A_1$ and $B=(X-a)B_1$, so
\[
          N=(X-a)^2(A_1^2-fB_1^2).
\]
This argument includes $B=0$ and handles branch points as well: a branch
point has only one ordinate and requires only one factor.

It follows that
\[
              \prod_a(X-a)^{m_a}\ \text{divides}\ N.
 \label{eq:weighted-roots}
\]
For completeness, this simultaneous divisibility uses no factorization
oracle.  Process the distinct $a$'s in their list order.  Suppose the
product $C$ of the already processed factors has been removed, so
$N=CQ$.  At the next coordinate $a$, every factor of $C$ has nonzero
value at $a$; hence $C(a)\ne0$.  The equality $N(a)=0$ gives $Q(a)=0$,
and synthetic division removes one $X-a$ from $Q$.  When $m_a=2$,
the known divisibility of $N$ by $(X-a)^2$, followed by cancellation of
one $X-a$, gives a second vanishing at $a$ and permits a second division.
All cancellations are by nonzero polynomials in $\mathbb Q[X]$.
This proves the displayed product divisibility by a finite induction.

The product is monic and has degree $\sum_a m_a$, which is exactly the
length of the original list.  Since $N\ne0$, division cannot increase
degree.  Therefore
\[
           \#\text{points}=\sum_a m_a\le\deg N\le L.
\]
Every operation just described has an explicit finite implementation.
Polynomial division by $X-a$ is synthetic division along the coefficient
list, and testing $A=0$ or $B=0$ is a finite coefficient check.
\end{proof}

\paragraph{Why counting only distinct $x$-coordinates loses information.}
A rough argument would count roots of $N$ and then multiply by two.  The
theorem avoids that loss because a fiber containing two common zeros
already consumes two factors of $N$.  This improvement is attained.
Take
\[
       f=X^3-X+1,\qquad A=X(X-1),\qquad B=0.
\]
Then $N=X^2(X-1)^2$ has degree four.  Its common zeros on the curve are
exactly
\[
                  (0,1),\ (0,-1),\ (1,1),\ (1,-1).
\]
Indeed $A(x)=0$ forces $x=0$ or $x=1$, and in either case $f(x)=1$.
Thus the bound four is exact; the rough $2\deg N$ bound would be eight.

\subsection{An explicit replacement for the divisor-based fiber contract}

Fix $L\ge d$ and put
\[
 a=\lfloor L/2\rfloor,\qquad b=\lfloor(L-d)/2\rfloor,\qquad
 n=a+b+2.
\]
Use the polynomial evaluation vector
\[
 \Phi_L(x,y)=(1,x,\ldots,x^a,\ y,xy,\ldots,x^b y)
                      \in\mathbb Q^n.
\]
The weight two assigned to $x$ and the weight $d$ assigned to $y$
are only a convenient organization of the norm degree in the proof;
no theory of divisors is being assumed.

\begin{corollary}[Proper subspaces have uniformly bounded fibers]
\label{cor:subspace-fiber}
For every proper rational linear subspace $W\subsetneq\mathbb Q^n$,
at most $L$ distinct rational points of $C_f$ have $\Phi_L(x,y)\in W$.
More precisely, every such finite list has length at most $L$.
\end{corollary}

\begin{proof}
Represent $W$ by a finite matrix of spanning vectors and the assertion
that its row rank is less than $n$.  Rational Gaussian elimination finds
a nonzero vector in its annihilator: select a free column, assign its
unknown the value one, assign the other free unknowns zero, and solve
backwards for the pivot unknowns.  This is a bounded exact procedure.
Write the resulting linear form on $\Phi_L$ as $A(x)+B(x)y$.
Its coefficient vector is nonzero, so $A,B$ are not both zero, and
its displayed degree bounds imply the hypotheses of
Theorem~\ref{thm:hyperelliptic-fiber}.  Every point mapping into $W$
is a zero of this form.  Apply that theorem.
\end{proof}

\begin{corollary}[A finite-tuple approximation theorem suffices]
\label{cor:hyperelliptic-cardinality}
Suppose a proved arithmetic contract says that for every finite list
of distinct points in a specified decidable subset $S\subseteq C_f$,
the corresponding vectors $\Phi_L$ can be covered by at most $M$ proper
rational subspaces.  Then every finite list of distinct points in $S$
has length at most $ML$.
\end{corollary}

\begin{proof}
Given the finite covering list, assign each point to the first subspace
containing its vector.  Membership is a finite rational rank test, and
coverage makes the bounded scan succeed.  Each of at most $M$ groups
contains at most $L$ points by Corollary~\ref{cor:subspace-fiber}.
Adding the group sizes gives the bound $ML$.  Equivalently, applying
this argument to a hypothetical list of $ML+1$ distinct points yields
an arithmetic contradiction.
\end{proof}

This explicitly discharges the zero-count half of the proposed Siegel
strategy for the odd-degree hyperelliptic evaluation systems above.
It does not prove the quantitative approximation theorem needed for the
covering contract, and hence does not prove general Siegel or Faltings.
The result supplies a cardinality bound once that contract is established;
it does not locate the points or give a coordinate cutoff.  The repair is
nevertheless substantive: one entire geometric dependency, previously
expressed through Riemann--Roch functions and zero/pole divisors, is
replaced here by a proved finite arithmetic lemma on this class.

\subsection{Finite algebra certificates as a general zero-count interface}

The preceding norm argument is particularly efficient.  A more general
interface works for arbitrary finite polynomial systems and can be reused
where a small norm formula is unavailable.

Let $F_1,\ldots,F_s\in\mathbb Q[X_1,\ldots,X_r]$.  A
\emph{spanning certificate of size $D$}, with $D\ge1$, consists of
polynomials $b_1=1,b_2,\ldots,b_D$, rational numbers $m_{kji}$, and
polynomials $u_{kj\ell}$ satisfying the finitely many identities
\[
 X_k b_j-\sum_{i=1}^D m_{kji}b_i
       =\sum_{\ell=1}^s u_{kj\ell}F_\ell
 \quad(1\le k\le r,\ 1\le j\le D).
 \label{eq:spanning-certificate}
\]
The checker expands and compares polynomial coefficients.  It checks
fixed algebraic identities, not the truth of arbitrary formal proofs.
Linear independence of the $b_i$ is not required; the certificate gives
an upper bound on dimension and can therefore be deliberately redundant.

\begin{theorem}[Certificate bound on rational zeros]
\label{thm:finite-algebra-zeros}
If a size-$D$ spanning certificate is valid, every finite list of distinct
common rational zeros of the $F_\ell$ has length at most $D$.
The checker, the reduction operations, and the finite contradiction
procedure used in the proof are primitive recursive.
\end{theorem}

\begin{proof}
The identities imply, by multiplying one variable at a time, that every
polynomial is congruent modulo $(F_1,\ldots,F_s)$ to a rational linear
combination of the $b_i$.  One can implement this reduction without a
Gr\"obner-basis algorithm: start with the coefficient vector of $1$,
apply the relevant matrix $m_k$ for each variable occurrence in a
monomial, and add the resulting vectors with the monomial coefficients.
For correctness, substitute a common zero $p$ in each certificate
identity.  Induction on the finite multiplication sequence shows that
the resulting coefficient vector $c$ satisfies
\[
                   H(p)=\sum_i c_i b_i(p)
\]
for the polynomial $H$ being reduced.  This equality is all that the
argument needs.

Suppose distinct common zeros $p_1,\ldots,p_t$ have been supplied.
For each ordered pair $i\ne j$, scan the $r$ coordinates to choose
the first $k(i,j)$ where $p_i$ and $p_j$ differ.  Their distinctness
makes the scan succeed.  Form
\[
 H_i(X)=\prod_{j\ne i}
 \frac{X_{k(i,j)}-(p_j)_{k(i,j)}}
      {(p_i)_{k(i,j)}-(p_j)_{k(i,j)}}.
\]
Every displayed denominator is a witnessed nonzero rational.  Direct
substitution gives $H_i(p_j)=1$ if $i=j$, and zero otherwise.
Reduce $H_i$ to a coefficient vector $c_i\in\mathbb Q^D$ by the
certificate matrices.  If a linear relation
$\sum_i\lambda_i c_i=0$ holds, evaluation at $p_j$ gives
\[
       0=\sum_i\lambda_i H_i(p_j)=\lambda_j.
\]
Thus the $t$ coefficient vectors are linearly independent.

Here the last dimension step is an explicit finite one.  If $t>D$,
perform rational row reduction on the $t$ by $D$ matrix with rows $c_i$,
carrying along the same row operations on the $t$ by $t$ identity matrix.
There are at most $D$ pivot rows, so at least one resulting coefficient
row is zero.  The corresponding row of the companion matrix is nonzero,
because every elementary row operation is invertible.  It gives a
nonzero rational relation among the $c_i$, contrary to the preceding
evaluation argument.  This proves $t\le D$.

All loops are bounded by the input lengths, numbers of variables,
exponents, and matrix dimensions.  Arithmetic may create large rational
coefficients, but addition, multiplication, normalization, and division
by the selected nonzero pivots are primitive recursive on rational codes.
No unbounded search for a basis, an algebraic closure, or a solution occurs.
\end{proof}

The finite common-zero theorem works uniformly in the encoded polynomial
system and in the certificate.  This is legitimate uniformity: polynomial
evaluation and coefficient comparison are specific primitive recursive
algorithms.  It is not a universal evaluator for arbitrary primitive
recursive programs.

\paragraph{A size-four certificate for the sharp example.}
For
\[
 F_1=y^2-x^3+x-1,\qquad F_2=x^2-x,
\]
take $b=(1,x,y,xy)$.  The elementary identities
\[
        x^2=x+F_2,\qquad
        y^2=1+F_1+(x+1)F_2
\]
and their multiples by $x$ or $y$ give all eight required border
identities.  Multiplication by $x$ sends the four basis entries, modulo
the two equations, to $(x,x,xy,xy)$; multiplication by $y$ sends them
to $(y,xy,1,x)$.  Therefore the certificate theorem gives at most four
rational solutions, and the four displayed points show that this bound
is attained.  The coefficient identities themselves can be checked by
an exact polynomial program without any analytic or geometric premise.

One should separate certificate \emph{verification} from certificate
\emph{construction}.  A system can have infinitely many zeros and then
admit no such certificate.  Even for a system known to be zero-dimensional,
calling an unspecified Gr\"obner-basis implementation does not by itself
prove a primitive recursive bound or an iPRA correctness theorem for its
construction.  Supplying a concrete certificate, or proving a bounded
constructor for the intended family, is the remaining input.

\subsection{An induction audit and Lean completion path}

Both zero-count arguments have particularly simple object-theory
invariants.  For Theorem~\ref{thm:hyperelliptic-fiber}, encode a finite
point list, its distinct-coordinate groups, a residual polynomial, and
the product of the factors already removed.  The loop invariant is the
conjunction of polynomial coefficient equalities
\[
 N=CQ,\qquad
 \deg N=\deg C+\deg Q,
\]
the accumulated root multiplicity, and the finite validity checks on the
processed groups.  It is used only while $N,Q$ are nonzero.  Nonzero
coefficient tests, polynomial degrees, list membership, and bounded
universal tests are all represented by primitive recursive Booleans.
Synthetic division carries its own coefficient identity as a second
finite invariant.  Consequently the induction formula contains no
unbounded quantifier over points or polynomials.  The arbitrary input
point list is a parameter.

For Theorem~\ref{thm:finite-algebra-zeros}, the reduction invariant is
evaluation equality for all entries of the supplied finite point list,
expressed by one bounded Boolean fold.  Row reduction carries the
coefficient matrix, the companion matrix, and its inverse; the equations
that these two square matrices multiply to the identity prove that the
companion row used at the end is nonzero.  Pivot selection is a bounded
scan with an explicit failure branch.  A failed pivot means the column
has no candidate, rather than an appeal to choice.  These are again
quantifier-free equalities on coded finite objects.

The formalization can be organized in the following order.
\begin{enumerate}
\item Define canonical rational codes, polynomial coefficient lists, finite
multivariate monomial lists, their primitive recursive evaluation
programs, and Boolean equality tests.  Prove the field and polynomial
identities used by synthetic division and matrix arithmetic in the
object theory.
\item Prove the synthetic-division identity and its specialization to a
root.  Prove nonvanishing and degree of products from leading coefficients.
Implement the distinct-coordinate grouping procedure and its
length-preservation identity.
\item Formalize the multiplicity-two fiber lemma: two distinct ordinates
annihilated by $A+B y$ force $A=B=0$ at that abscissa.  Combine its
explicit factors with the root-removal invariant to obtain the
hyperelliptic bound.
\item Define the spanning-certificate Boolean by coefficientwise checks.
Prove the reduction invariant, construct the finite separator vectors,
and prove their evaluation matrix is the identity.
\item Formalize bounded Gaussian elimination, including the tracked
inverse for row operations.  Use it for both a proper-subspace
annihilator and the finite-algebra dimension contradiction.
\item Assemble actual embedded derivations of the zero-count formulas.
Only after these are constructed should the existing ConciseIPRA
interpretation machinery transport them to another presentation of the
theory.
\end{enumerate}

For example, the native specification of the final fiber statement has
the following shape, where every predicate in the hypotheses is a
decidable finite-code predicate and \texttt{points.length} refers to an
encoded list length:
\begin{verbatim}
def HyperellipticFiberBound : Prop :=
  forall f A B points,
    OddDegree f -> (Nonzero A || Nonzero B) = true ->
    PairwiseDistinct points = true ->
    AllCurveAndLinearZeros f A B points = true ->
    points.length <= degree (A*A - f*B*B)

def FiniteAlgebraZeroBound : Prop :=
  forall equations certificate points,
    CheckSpanning equations certificate = true ->
    PairwiseDistinct points = true ->
    AllCommonZeros equations points = true ->
    points.length <= certificate.dimension
\end{verbatim}
This pseudocode specifies the mathematical types; it is not presented as
a compiled Lean module.  In an actual implementation, the polynomial
operators, code projections, and Boolean names must be replaced by the
chosen concrete definitions.  The repository endpoint is a corresponding
object-language formula together with a term of type
\texttt{Prf PRAAx [] formula}, not merely a proof of the native
\texttt{Prop}.  Neither a field in a Lean structure nor an uninstantiated
contract should be counted as that derivation.

